\documentclass[12pt]{exam}
\usepackage[margin=1.5in]{geometry}

\usepackage{
  % AH> I don't know how much of these are used, and for what.
  %     (needs untangling.) 
  enumitem,
  xspace,
  mathtools,
  multicol,
  tikz-cd,
  parskip,
  mathrsfs,
  stmaryrd,
  tikz,
  float,
  titlecaps,
  soul,
  upgreek,
  caption,
  graphicx,
  array,
  subcaption,
  ../sty/mathwidth,
%  ./sty/sectsty,
%  ./sty/commands,
  ../sty/infer,
  ../sty/commands}


% Minted (needs Python)
% \usepackage{minted}

% tikz
\usepackage{tikz}
\usepackage{tikz-cd}


\title{Select Solutions to Introduction to Homotopy Type Theory, Egbert Rijke}

\author{Computational Logic Center}


\begin{document}

\maketitle

\section*{1.1}
\label{sec:ChapterOne}

\newcommand\Jeq{\ensuremath{\dot{=}}}
\newcommand\Eq[2]{\ensuremath{#1\, \Jeq\, #2}}
\newcommand\EqJ[4]{\ensuremath{#1 \vdash \Eq {#2} {#3} : #4}}
\newcommand\Type[2]{\ensuremath{#1 \vdash #2 \; \text{type}}}
\newcommand\TypeJ[3]{\ensuremath{#1 \vdash #2 : #3}}

\begin{enumerate}[label=(\alph*)]
  \item Give a derivation for the following \textbf{element conversion rule}:
  \[ 
    \ib{\irule
      {\Type{\Gamma}{\Eq {A} {A'}}}
      {\TypeJ \Gamma a A};
      {\TypeJ {\Gamma} {a} {A'}}}
  \]
  \item Give a derivation for the following \textbf{congruence rule} for element conversion:
  \[ 
    \ib{\irule
      {\Type \Gamma {\Eq {A} {A'}}}
      {\EqJ \Gamma {a} {b} {A}}; 
      {\EqJ \Gamma {a} {b} {A'}}
    }
  \]
\end{enumerate}

\subsubsection*{Solution (a)}

\begin{gather*}
{\small
    \ib{\irule[subst]
      {\Gamma \vdash a : A}
      {\irule[var-conv]
        {\irule[sym]{\Type \Gamma {\Eq {A} {A'}}}; {\Type \Gamma {\Eq {A'} {A}}}}
        {\irule[var]{};{\Gamma, x : A' \vdash x : A'}}; 
        {\Gamma , x : A \vdash x : A'}};
      {\TypeJ {\Gamma} {x[a/x]} {A'[a/x]}}}} 
\end{gather*}

\subsubsection*{Solution (b)}

\begin{gather*}
    \ib{\irule[cong-subst]
      {\EqJ \Gamma a b A}
      {\irule[above]{};{\TypeJ {\Gamma , x : A} x A'}}; 
      {\EqJ \Gamma {x[a/x]} {x[b/x]} {A'[a/x]}}
    }  
\end{gather*}

\section*{2.1}

The $\eta$-rule is often seen as a judgmental extensionality principle. Use the $\eta$-rule to show that if $f$ and $g$ take equal values, then they must be equal, i.e., give a derivation of the rule:

\begin{gather*}
\ib{\irule
  {\begin{array}{@{}c@{}}
             {\TypeJ \Gamma f {\Pi_{x : A}B(x)}}
             \\
             {\TypeJ \Gamma g {\Pi_{x : A}B(x)}}
             \\
             {\EqJ {\Gamma , x : A} {f(x)}{g(x)} {B(x)}}
           \end{array}};
  {\EqJ \Gamma f g {\Pi_{x : A}B(x)}}}
\end{gather*}

\subsection*{Solution}
Permit me some creative liberty in chaining two invocations of the transitivity rule.

{\footnotesize
\begin{gather*}
\ib{\irule[trans]
    {\irule[\ensuremath{\eta}]
      {\TypeJ \Gamma f {\Pi_{x : A}B(x)}};
      {\EqJ \Gamma {f} {\lambda x. f(x)} {\Pi_{x : A}B(x)}}}
    {\irule[$\lambda$-eq]{\EqJ {\Gamma , x : A} {f(x)} {g(x)} {B(x)}};{\EqJ \Gamma {\lambda x. f(x)}{\lambda x. g(x)} {B(x)}}}
    {\irule[\ensuremath{\eta}]
      {\TypeJ \Gamma g {\Pi_{x : A}B(x)}};
      {\EqJ \Gamma {\lambda x. g(x)} {g} {\Pi_{x : A}B(x)}}};
  {\EqJ \Gamma f g {\Pi_{x : A}B(x)}}}
\end{gather*}}


\bibliographystyle{plainnat}
\bibliography{template.bib}
\end{document}


%%% Local Variables:
%%% mode: latex
%%% eval: (visual-line-mode 1)
%%% eval: (auto-fill-mode 0)
%%% eval: (tex-source-correlate-mode 1)
%%% eval: (flyspell-mode 1)
%%% TeX-command-extra-options: "--synctex=1"
%%% TeX-master: t
%%% End:
